\documentclass{lomonosov}

\begin{thesis}

\Title{Верифкация ИНС на основе тензорного парллелизма}{{Воробьев\,С.\,Ю.}}

\Author{Воробьев~Сергей~Юрьевич}{Студент}{Факультет ВМК МГУ имени М.\,В.\,Ломоносова}{Москва}{Россия}{sergei.vorobyov01@yandex.ru}{Ильюшин~Евгений~Альбинович}

Формальная верификация нейронных сетей~--- задача математического доказательства того, что сеть удовлетворяет заданным спецификациям (например, робастность к состязательным возмущениям входных данных) для всех допустимых входов. Для сетей с функцией активации \ENGLISH{ReLU} задача является NP-полной~[1], что обуславливает экспоненциальный рост вычислительных затрат с увеличением размера модели. Ведущим инструментом верификации является \ENGLISH{Alpha-Beta-CROWN}~--- победитель международного соревнования \ENGLISH{VNN-COMP} в 2021--2024 годах~[2]. Инструмент использует метод распространения линейных границ (\ENGLISH{Linear Relaxation based Perturbation Analysis, LiRPA}) с оптимизируемыми параметрами релаксации $\alpha$ и двойственными множителями Лагранжа $\beta$, реализованный в виде операций над тензорами на GPU в библиотеке \texttt{auto\_LiRPA}~[3].

С ростом масштабов верифицируемых моделей (\ENGLISH{Transformers}, большие языковые модели) возникает проблема ограниченной памяти и вычислительных ресурсов одного GPU. В области обучения нейронных сетей аналогичная проблема решается методами параллелизма, в частности \ENGLISH{Tensor Parallelism} (TP), впервые предложенным в работе~[4]. Цель настоящего исследования~--- адаптация TP к задаче верификации нейронных сетей методом распространения границ.

Основная идея состоит в распределении весовых матриц линейных слоёв между несколькими GPU с параллельным выполнением обратного распространения границ (\ENGLISH{CROWN backward pass}~[5]). Реализованы два типа параллельных слоёв: \ENGLISH{Column Parallel}, в котором матрица весов разделяется по выходной размерности, а результаты агрегируются операцией \ENGLISH{AllReduce}; и \ENGLISH{Row Parallel}, в котором матрица разделяется по входной размерности, и коммуникация при обратном проходе не требуется. Чередование этих слоёв образует схему, минимизирующую объём межпроцессных коммуникаций.

Реализация выполнена в виде расширений \ENGLISH{auto\_LiRPA} для распределённых вычислений через коммуникации между GPU. При использовании $N$ GPU потребление памяти весов и матриц границ на каждом ускорителе снижается в $N$ раз: для MLP-блока с размерностью $d = 4096$ это соответствует сокращению с 512 до 256~МБ при $N = 2$. Замеры на конфигурации из двух \ENGLISH{NVIDIA A100}, соединённых \ENGLISH{NVLink}, показали ускорение вычисления границ \ENGLISH{CROWN} в ${\sim}1{,}97$ раза при накладных расходах на коммуникации \ENGLISH{AllReduce} менее 2\% от общего времени. Полученные результаты демонстрируют принципиальную возможность масштабирования формальной верификации нейронных сетей на несколько GPU, что открывает путь к верификации более крупных архитектур.

\begin{references}
\Source \ENGLISH{Ro\-bust\-ness Ve\-ri\-fi\-ca\-tion in Neu\-ral Net\-works~// arXiv pre\-print arXiv:2403.13441. 2024.}

\Source \ENGLISH{The Fifth In\-ter\-na\-tion\-al Ve\-ri\-fi\-ca\-tion of Neu\-ral Net\-works Com\-pe\-ti\-tion (VNN-COMP 2024): Sum\-ma\-ry and Re\-sults~// arXiv pre\-print arXiv:2412.19985. 2024.}

\Source \ENGLISH{Xu\,K., Shi\,Z., Zhang\,H. et al. Au\-to\-mat\-ic Per\-tur\-ba\-tion Ana\-ly\-sis for Scal\-able Cer\-ti\-fied Ro\-bust\-ness and Be\-yond~// Ad\-vanc\-es in Neu\-ral In\-for\-ma\-tion Pro\-ces\-sing Sys\-tems~33 (NeurIPS). 2020.}

\Source \ENGLISH{Shoey\-bi\,M., Pat\-wa\-ry\,M., Pu\-ri\,R. et al. Me\-ga\-tron-LM: Train\-ing Mul\-ti-Bil\-lion Pa\-ram\-e\-ter Lan\-guage Mo\-dels Us\-ing Mo\-del Pa\-ral\-lel\-ism~// arXiv pre\-print arXiv:1909.08053. 2020.}

\Source \ENGLISH{Zhang\,H., Weng\,T.-W., Chen\,P.-Y. et al. Ef\-fi\-cient Neu\-ral Net\-work Ro\-bust\-ness Cer\-ti\-fi\-ca\-tion with Gen\-er\-al Ac\-ti\-va\-tion Func\-tions~// Ad\-vanc\-es in Neu\-ral In\-for\-ma\-tion Pro\-ces\-sing Sys\-tems~31 (NeurIPS). 2018.}
\end{references}

\end{thesis}
